\documentclass[11pt]{article}

\usepackage[a4paper,margin=1in]{geometry}
\usepackage{amsmath,amssymb}
\usepackage{booktabs}
\usepackage{hyperref}
\usepackage{listings}
\usepackage{xcolor}
\usepackage{enumitem}
\usepackage{xurl}   % best at breaking long strings

\lstset{
  basicstyle=\ttfamily\small,
  breaklines=true,
  frame=single,
  columns=fullflexible
}

\title{Containment-Aware Agreement for GeoJSON Annotations with Split--Merge Robust Evaluation}
\author{}
\date{}

\begin{document}
\maketitle

\begin{abstract}
This report presents a containment-aware evaluation framework for inter-annotator agreement on region annotations stored in GeoJSON, motivated by histopathology workflows in which annotators may differ substantially in localization strictness and in whether they split or merge regions.
Conventional Intersection-over-Union (IoU) can underestimate agreement when one annotator draws strict (tight) regions and another draws loose (coarse) regions that nonetheless contain the relevant tissue structures.
The proposed framework replaces IoU with a \emph{Containment-Aware Similarity} (CAS) score and provides three complementary evaluation modes: (i) one-to-one instance matching, (ii) union-per-class area agreement, and (iii) connected-component grouping for split--merge robust count-based scoring.
A directory-wise implementation evaluates all matching subdirectories between two directory roots and optionally exports merged GeoJSON files with match metadata for qualitative inspection.
\end{abstract}

\section{Introduction}
Inter-annotator agreement for detection-style annotations is often evaluated by matching predicted and reference instances using IoU thresholds and reporting precision/recall/F1 or average precision.
However, IoU assumes that annotators delineate regions with comparable tightness.
In histopathology and related imaging tasks, one annotator may annotate only the minimal region corresponding to a structure of interest, while another may annotate a larger free-hand region that includes the same structure plus surrounding tissue.
Although these annotations may be semantically aligned, IoU decreases due to a large union area relative to the intersection.

A second common failure mode arises from \emph{split--merge} differences: one annotator may represent a single semantic region using multiple small boxes, while another uses a single large box.
Strict one-to-one matching then produces artificially low agreement by forcing at most one small region to match the large region.
This report describes a framework that addresses both issues by using CAS and by offering evaluation modes that are robust to split--merge variability.

\section{Data Model and Problem Setup}
Two annotation sources are compared for each case (folder):
\begin{itemize}[leftmargin=*]
  \item \textbf{Annotator A:} treated as the reference set for evaluation purposes.
  \item \textbf{Annotator B:} treated as the compared set for evaluation purposes.
\end{itemize}

Annotations are stored as GeoJSON \texttt{FeatureCollection}s where each feature contains:
\begin{itemize}[leftmargin=*]
  \item a geometry (Polygon or MultiPolygon),
  \item a class label stored typically at \texttt{properties.classification.name}.
\end{itemize}

The evaluation is performed \emph{per class} and aggregated \emph{per folder}, enabling case-level summaries across many slides or regions.

\section{Containment-Aware Similarity (CAS)}
Let $P$ denote a region from annotator B and $G$ a region from annotator A.
Denote their areas by:
\[
A_P = |P|,\quad A_G = |G|,\quad A_I = |P \cap G|.
\]

\subsection{Overlap Coefficient}
A containment-friendly overlap measure is the overlap coefficient:
\begin{equation}
\mathrm{OC}(P,G) = \frac{A_I}{\min(A_P, A_G)}.
\end{equation}
This coefficient equals $1$ when the smaller region is fully contained in the larger region and therefore does not penalize additional area in loose annotations.

\subsection{Symmetric Area Ratio}
A symmetric size mismatch measure is:
\begin{equation}
\mathrm{AR}(P,G) = \frac{\max(A_P, A_G)}{\min(A_P, A_G)} \ge 1.
\end{equation}

\subsection{CAS Definition}
CAS is defined as the overlap coefficient multiplied by a smooth size-mismatch penalty:
\begin{equation}
\mathrm{CAS}(P,G) = \mathrm{OC}(P,G)\cdot \exp\left(-\frac{(\ln \mathrm{AR}(P,G))^2}{2\sigma^2}\right),
\end{equation}
where $\sigma>0$ controls tolerance to size mismatch; larger $\sigma$ yields weaker penalization.
When $\sigma \rightarrow \infty$ (or the penalty is disabled), CAS reduces to $\mathrm{OC}$.

\paragraph{Interpretation.}
CAS preserves the containment property of $\mathrm{OC}$ while discouraging degenerate cases in which extremely large regions trivially contain many smaller regions.

\section{Agreement Evaluation Modes}
The framework provides three evaluation modes, each addressing different annotation behaviors.

\subsection{Mode 1: One-to-One Instance Matching (\texttt{one\_to\_one})}
For each class $c$, the sets of instances $\{G_i\}_{i=1}^{n}$ (annotator A) and $\{P_j\}_{j=1}^{m}$ (annotator B) are compared.

\paragraph{Candidate scoring.}
CAS is computed for all candidate pairs $(P_j,G_i)$.
For efficiency at scale, candidate generation can be accelerated using spatial indexing (e.g., R-tree queries).

\paragraph{One-to-one assignment.}
A one-to-one matching is computed using either:
\begin{itemize}[leftmargin=*]
  \item greedy maximum-CAS selection over pairs with $\mathrm{CAS}\ge\tau$, or
  \item Hungarian maximum-weight assignment for small problems, with dummy nodes enabling unmatched items.
\end{itemize}

\paragraph{Counts and metrics.}
At a CAS threshold $\tau$:
\begin{itemize}[leftmargin=*]
  \item $\mathrm{TP}$ equals the number of matched pairs,
  \item $\mathrm{FP}$ equals unmatched B instances,
  \item $\mathrm{FN}$ equals unmatched A instances.
\end{itemize}
Precision, recall, and F1 are computed in the usual manner.

\paragraph{Limitation.}
One-to-one matching can underestimate agreement when split--merge behavior is present (e.g., one large instance corresponding to multiple smaller instances).

\subsection{Mode 2: Union-Per-Class Area Agreement (\texttt{union\_per\_class})}
This mode removes instance matching entirely and is robust to split--merge variability.
For each class $c$, union geometries are formed:
\[
A_c = \bigcup_i G_i,\qquad B_c = \bigcup_j P_j.
\]
Area-based precision and recall are defined by:
\begin{equation}
\mathrm{Precision}_c = \frac{|A_c\cap B_c|}{|B_c|},\qquad
\mathrm{Recall}_c = \frac{|A_c\cap B_c|}{|A_c|}.
\end{equation}
F1 is computed from these quantities.

\paragraph{Diagnostics.}
In addition, $\mathrm{CAS}(A_c,B_c)$, $\mathrm{OC}(A_c,B_c)$, and $\mathrm{AR}(A_c,B_c)$ are reported as union-level diagnostics, along with areas $|A_c|$, $|B_c|$, and $|A_c\cap B_c|$.

\paragraph{Property.}
Union-per-class scoring is invariant to instance splitting/merging within a class, provided the unioned regions remain the same.

\subsection{Mode 3: Connected-Component Split--Merge Grouping (\texttt{components})}
This mode provides count-based agreement while remaining robust to split--merge variability.
For each class $c$, a bipartite graph is constructed:
\begin{itemize}[leftmargin=*]
  \item nodes correspond to individual A instances and B instances,
  \item an edge connects $G_i$ and $P_j$ if $\mathrm{CAS}(P_j,G_i)\ge\tau$.
\end{itemize}
Connected components in this graph represent \emph{groups} of instances believed to refer to the same semantic region under split--merge behavior.

\paragraph{Component counting.}
Let a component contain $n_A$ A-instances and $n_B$ B-instances.
It is categorized as:
\begin{itemize}[leftmargin=*]
  \item \textbf{TP component} if $n_A\ge 1$ and $n_B\ge 1$,
  \item \textbf{FN component} if $n_A\ge 1$ and $n_B=0$,
  \item \textbf{FP component} if $n_A=0$ and $n_B\ge 1$.
\end{itemize}
Component-level precision/recall/F1 are then computed using TP/FP/FN component counts.

\paragraph{Component quality diagnostics.}
For each component, union geometries are computed:
\[
A_k = \bigcup_{i\in k} G_i,\qquad B_k = \bigcup_{j\in k} P_j,
\]
and union-level CAS/OC/AR and areas are reported as diagnostics.

\paragraph{Advantages.}
Unlike one-to-one matching, this mode does not penalize a single large annotation being represented by multiple smaller annotations, yet it avoids the inflation risk of unconstrained many-to-many matching by scoring at the component level.

\section{Directory-Level Evaluation Protocol}
The implementation evaluates agreement across entire directory trees.
Let $A_{\mathrm{root}}$ and $B_{\mathrm{root}}$ be directory roots.
The procedure is:
\begin{enumerate}[leftmargin=*]
  \item Recursively enumerate all \texttt{*.geojson} files under each root.
  \item Group files by their \emph{relative parent directory} (relative to the root).
  \item Identify keys present in both roots and treat each key as a matched case.
  \item For each matched case, load all features from the GeoJSON files under that folder and evaluate agreement per class using the selected mode.
\end{enumerate}

\section{Outputs}
For each matched folder key, the script writes:
\begin{itemize}[leftmargin=*]
  \item \texttt{per\_class\_metrics.csv}: per-class precision/recall/F1; TP/FP/FN semantics depend on mode,
  \item \texttt{summary.json}: micro/macro summaries and configuration parameters,
  \item \texttt{matches.csv}: detailed match records.
\end{itemize}
A global \texttt{cases\_summary.csv} and \texttt{overall.txt} are also written.

\subsection{Mode-Dependent Semantics of \texttt{matches.csv}}
\begin{itemize}[leftmargin=*]
  \item \textbf{one\_to\_one:} one row per matched pair with CAS/OC/AR and areas.
  \item \textbf{union\_per\_class:} one row per class with union-level CAS/OC/AR and union areas.
  \item \textbf{components:} one row per connected component with membership lists and union-level diagnostics.
\end{itemize}

\subsection{Optional GeoJSON Exports}
When enabled, the script exports \texttt{combined\_AB.geojson} containing all features from both sources with source-tagged labels (e.g., ``A $|$ ClassName'') and an attached \texttt{match} property carrying match metadata.
Additionally:
\begin{itemize}[leftmargin=*]
  \item \texttt{matched\_AB.geojson} is written in \textbf{one\_to\_one} mode containing only matched pairs.
  \item \texttt{matched\_AB.geojson} is written in \textbf{components} mode containing only TP components.
  \item \texttt{matched\_AB.geojson} is not meaningful in \textbf{union\_per\_class} mode and is therefore omitted.
\end{itemize}

\section{Usage Examples}
One-to-one matching:
\begin{lstlisting}
python cas_agreement_dir_eval_v3.py \
  --a_root /path/to/A \
  --b_root /path/to/B \
  --out cas_out \
  --mode one_to_one \
  --cas_thr 0.5 \
  --sigma 1.0 \
  --matching greedy \
  --write_geojson
\end{lstlisting}

Union-per-class:
\begin{lstlisting}
python cas_agreement_dir_eval_v3.py \
  --a_root /path/to/A \
  --b_root /path/to/B \
  --out cas_out \
  --mode union_per_class \
  --sigma 1.0 \
  --write_geojson
\end{lstlisting}

Connected components:
\begin{lstlisting}
python cas_agreement_dir_eval_v3.py \
  --a_root /path/to/A \
  --b_root /path/to/B \
  --out cas_out \
  --mode components \
  --cas_thr 0.5 \
  --sigma 1.0 \
  --write_geojson
\end{lstlisting}

\section{Discussion and Limitations}
Several considerations apply:
\begin{itemize}[leftmargin=*]
  \item \textbf{Choice of $\tau$ and $\sigma$:} the CAS threshold $\tau$ controls what constitutes a link/match, and $\sigma$ controls tolerance to size mismatch. Both should be selected based on expected annotation variability.
  \item \textbf{Degenerate containment:} overlap-coefficient-based measures can be high for overly large regions; the size penalty mitigates this, but extreme cases should be inspected.
  \item \textbf{Class taxonomy consistency:} consistent class naming between sources is assumed; label normalization may be required.
  \item \textbf{No confidence modeling:} as annotations are manual, standard confidence-based AP/mAP is not applicable. The framework reports thresholded agreement metrics and optional diagnostic summaries.
\end{itemize}

\section{Conclusion}
This report described a containment-aware agreement framework for GeoJSON region annotations that is robust to differing localization strictness and split--merge annotation behavior.
By using CAS in place of IoU and providing three evaluation modes---one-to-one instance matching, union-per-class area agreement, and connected-component grouping---the framework supports both instance-level auditing and robust case-level agreement estimation.
The accompanying directory-wise implementation enables scalable evaluation and produces merged GeoJSON exports for qualitative validation.

\end{document}
